% Generated draft: requires mathematical and source review.
\documentclass[11pt]{article}
\usepackage[margin=25mm]{geometry}
\usepackage{fontspec}
\setmainfont{Latin Modern Roman}
\usepackage{amsmath,amssymb,mathtools,unicode-math}
\setmathfont{Latin Modern Math}
\usepackage{enumitem,needspace,xurl,hyperref}
\setlength{\parindent}{0pt}
\setlength{\parskip}{5pt}
\setlength{\emergencystretch}{3em}
\newcommand{\N}{\mathbb{N}}
\newcommand{\Z}{\mathbb{Z}}
\newcommand{\Q}{\mathbb{Q}}
\newcommand{\R}{\mathbb{R}}
\newcommand{\C}{\mathbb{C}}
\newcommand{\F}{\mathbb{F}}
\newcommand{\A}{\mathbb{A}}
\newcommand{\PP}{\mathbb{P}}
\newcommand{\E}{\mathbb{E}}
\newcommand{\eps}{\varepsilon}
\newcommand{\dd}{\,\mathrm{d}}
\begin{document}
\section*{UnsolvedMath ID: 2307042}
\textbf{Source problem code:} catalogue entry
\subsection{Definitions and mathematical statement}
Let $D\subset\C$ be a Green domain (a domain admitting a Green function) and fix $z_0\in D$. Define the Harnack function
\[
H_{z_0}(z)=\sup\{h(z):h>0\text{ harmonic on }D,\ h(z_0)\leq1\}.
\]
The source repeats part of this defining sentence. Its intended Martin-kernel normalization appears to be $K_\zeta(z_0)=1$ for boundary parameters $\zeta\in\Delta_1$. Under the appropriate Martin representation, the natural displayed identity is
\[
H_{z_0}(z)=\sup_{\zeta\in\Delta_1}K_\zeta(z),\qquad z\in D.
\]
However, the supplied text instead has $K(\zeta_0)=1$ and $H_\zeta$ in the supremum, neither defined, and calls $\Delta_1$ the Martin boundary without clarifying whether it means the minimal boundary. These are explicit notation issues requiring verification.

For $D=\mathbb D$ and $z_0=0$, the source identifies $\Delta_1$ with $\mathbb T=\{|\zeta|=1\}$ and $K_\zeta$ with the normalized Poisson kernel. In this standard case,
\[
K_\zeta(z)=\frac{1-|z|^2}{|\zeta-z|^2},\qquad
H_0(z)=\frac{1+|z|}{1-|z|},\qquad
K_\zeta(r\zeta)=H_0(r\zeta)\quad(0\leq r<1).
\]
Thus a single boundary kernel realizes the supremum all along its associated radius. The source says conformal mapping gives the corresponding simply connected property: each $K_\zeta$ touches $H_{z_0}$ along a Green line issuing from $z_0$.

Arsove and Johnson ask whether this continues to hold in multiply connected Green domains. A natural exact equality reading is: for each relevant $\zeta\in\Delta_1$, is there a Green line $\ell_\zeta$ from $z_0$ such that $K_\zeta(z)=H_{z_0}(z)$ for every $z\in\ell_\zeta$? Green lines are commonly trajectories orthogonal to level curves of the Green function with pole $z_0$, but the record does not specify their continuation at critical points or their boundary endpoints. Nor does it formally define ``touches along'' or clarify whether an entire line or only a segment is required. Equality only at $z_0$ would be automatic from normalization and is not substituted for the stated along-a-line property. The supplied summary reports openness as of 2018 and no definitive resolution located.
\subsection{Short English statement}
In the disc, each normalized Poisson kernel equals the extremal Harnack function along one radius; the source reports an analogous Green-line statement in simply connected domains. The question is whether each Martin kernel has such a line in multiply connected domains, with boundary and line conventions needing clarification.
\subsection{Sources}
\begin{itemize}
\item Local UnsolvedMath registry, permanent ID 2307042, exact source code catalogue entry, ``Research Problems in Function Theory --- Problem 7.42''; status open, classification OPEN-TRIAGE.
\item Hayman, \emph{Research Problems in Function Theory} (2018), Problem 7.42, attributed to M. Arsove and G. Johnson, \url{https://arxiv.org/abs/1809.07200}. The record identifies extraction from public arXiv source TeX, accessed July 29, 2026; the 2018 update was retained for status triage.
\item The supplied statement gives the disc and simply connected comparison without separate bibliographic references; its Martin notation contains the inconsistencies recorded above.
\item The supplied research summary reports openness as of the 2018 edition and no definitive resolution located, with no additional authoritative reference or precise search date.
\end{itemize}
\end{document}
